\SourceLeaf{350}%
wesentlich verschiedene Relationen zu unterscheiden sind, theils nemlich solche, die eine bestimmte Form der Fläche im Raume voraussetzen, theils solche, welche von den verschiedenen Formen, die die Fläche annehmen kann, unabhängig sind. Die letztern sind es, wovon hier die Rede ist; nach dem, was vorhin bemerkt ist, gehört dazu das Krümmungsmaass; man sieht aber leicht, dass eben dahin die Betrachtung der auf der Fläche construirten Figuren, ihrer Winkel, ihres Flächeninhalts und ihrer Totalkrümmung, die Verbindung der Punkte durch kürzeste Linien u.\,dgl. gehört. Alle solche Untersuchungen müssen davon ausgehen, dass die Natur der krummen Fläche an sich durch den Ausdruck eines unbestimmten Linearelements in der Form \(\BareRoot(E\mathrm{d}p^2+2F\mathrm{d}p\mathord{.}\mathrm{d}q+G\mathrm{d}q^2)\) gegeben ist. Der Verf. hat gegenwärtiger Abhandlung einen Theil seiner seit mehreren Jahren auf diesem Felde angestellten Untersuchungen einverleibt, indem er sich auf solche einschränkte, die von dem ersten Eintritt nicht zu entfernt liegen und zum Theil als allgemeine Hülfsmittel zu vielfachen weitern Untersuchungen dienen können. Bei unsrer Anzeige müssen wir uns noch mehr beschränken, und uns begnügen, nur einiges als Probe anzuführen. Als solche mögen folgende Lehrsätze dienen.

Wenn auf einer krummen Fläche von Einem Anfangspunkte ein System unendlich vieler kürzester Linien von gleicher Länge ausläuft, so schneidet die durch ihre Endpunkte gehende Linie jede derselben unter rechten Winkeln. Wenn an jedem Punkte einer beliebigen Linie auf einer krummen Fläche kürzeste Linien von gleicher Länge senkrecht gegen jene Linie gezogen sind, so sind diese alle auch senkrecht gegen diejenige Linie, welche ihre andern Endpunkte verbindet. Diese beiden Lehrsätze, wovon der zweite als eine Generalisirung des ersten betrachtet werden kann, werden sowohl analytisch, als durch einfache geometrische Betrachtungen bewiesen. \textit{Der Überschuss der Summe der Winkel eines aus kürzesten Linien gebildeten Dreiecks über zwei Rechte ist der Totalkrümmung des Dreiecks gleich.} Es wird hiebei angenommen, dass für die Winkel derjenige. dem ein dem Halbmesser gleicher Bogen entspricht, \((57^0\ 17'\ 45'')\), und für die ganze Krümmung, als Stück der Fläche der Hülfskugel, der Inhalt eines Quadrats, dessen Seite der Halbmesser der Hülfskugel ist, als Einheit zum Grunde liegt. Offenbar kann man diess wichtige Theorem auch so ausdrücken: der Überschuss der Winkel eines aus kürzesten Linien gebildeten Dreiecks über zwei Rechte verhält sich zu acht Rechten, wie das Stück der Oberfläche der Hülfsku%
